% ============================================================
% kbs_method.tex — KBS journal Method (KE voice).
% Section label: sec:method (referenced by the Introduction draft).
% Builds on conference Method (paper/main.tex §III); journal extensions
% (corpus scale, multi-model, panel) are marked \TODO{E?}.
% Bibliography: kbs_ke_references.bib + paper/references.bib.
% ============================================================

\section{Method}
\label{sec:method}

We formalize the paper as a knowledge-engineering study. This section
defines the knowledge being acquired (Section~\ref{sec:method:knowledge}),
the acquisition component (Section~\ref{sec:method:acquisition}), the
knowledge-based system that consumes the knowledge
(Section~\ref{sec:method:system}), and the measurement apparatus used to
validate and evaluate it (Section~\ref{sec:method:benchmark}).

\subsection{Knowledge being acquired}
\label{sec:method:knowledge}

The unit of knowledge is the \emph{catalog record}: a description of a work
that pairs an identity (title, author, year) with subject knowledge about
the work. Subject knowledge takes two canonical representations, and we
study both.

\textbf{Representation 1 --- subject headings (thesaurus).} Subject headings
(e.g., Library of Congress Subject Headings) are an open, pre-coordinated
thesaurus: a large, human-curated vocabulary of heading strings that may
combine facets with subdivision markers (e.g., ``Land tenure -- England --
History -- To 1500''). New headings are legitimately added over time, and
different catalogers may assign different but equally valid headings for the
same work. This makes agreement-based evaluation subtle: two professional
sources do not agree exactly on every heading
\TODO{E2: quantify from panel}.

\textbf{Representation 2 --- classification (taxonomy).} Classification
(e.g., Dewey Decimal Classification) is a closed, hierarchical taxonomy of
numeric classes with a defined top-level structure and a single canonical
placement for a work's primary subject. It is more constrained than the
thesaurus and, as we show in Section~\ref{sec:util}, behaves differently
when machine-acquired.

A record may carry neither representation (no subject knowledge), one, or
both. The central quantity of this paper is \emph{knowledge-base
completeness}: the fraction of records in a collection that carry subject
knowledge in a given representation. Sparse records---those with no
headings---are invisible to topical knowledge-based retrieval no matter how
good the retrieval machinery is; this is the gap the acquisition component
is asked to fill.

\subsection{Acquisition component}
\label{sec:method:acquisition}

The acquisition component is an LLM configured to produce subject knowledge
from a sparse record. It receives only what a minimal record carries
---title, author, year---and must propose, per record: (i) subject headings
in LCSH form and (ii) a DDC number. This is the retrospective-conversion
setting: libraries hold millions of records with no subject knowledge, and
the question is whether an LLM can acquire it. The model is prompted without
the gold answer and, in the main condition, without the work's full text;
a full-text-conditioned variant is studied as an acquisition-lever ablation
in Section~\ref{sec:valid} \TODO{E3}. We report results for
\texttt{deepseek-chat} as the primary acquisition model, with
multi-model replication \TODO{E3} to bound model dependence.

\subsection{Consuming knowledge-based system}
\label{sec:method:system}

The consuming system is a metadata-aware retrieval index over the catalog:
the benchmark for whether acquired subject knowledge is \emph{useful} is
whether indexing it improves retrieval quality. We deliberately keep the
retrieval machinery standard---a field-weighted BM25F-style ranker
paired with a dense twin, fused by reciprocal-rank fusion---and claim no
novelty for it~\cite{bm25f, prf_bm25}; the contribution is the knowledge
analysis, not a new retriever. Two indexes are built over each collection:
full-text chunks (dense embeddings plus BM25) and \emph{catalog records} as
first-class retrieval units (field-boosted BM25 over title, author, and
subject-heading fields, with subject weight $\times2$, plus dense embeddings
of the field text). Given a query, records are ranked by RRF over the dense
and lexical record rankings; full-text chunks are retrieved in parallel and
used only as generation evidence, not in the record ranking. We evaluate
four configurations---\texttt{dense}, \texttt{bm25}, \texttt{hybrid}
(metadata-blind chunk retrieval), and \texttt{meta} (the record-aware
configuration)---so that the contribution of the subject-knowledge fields is
isolated by subtraction.

The representation-dependence question is operationalized as a
\emph{field ablation}: we rebuild the record index varying only which
knowledge fields enter the record text (authors-only baseline; +title;
+subject headings; +classification), holding retrieval, generation, and the
chunk index fixed, and score each configuration on the full question pool
(Section~\ref{sec:util}). This is what lets us attribute downstream utility
to a specific knowledge representation rather than to ``metadata'' as an
undifferentiated whole.

\subsection{Measurement apparatus}
\label{sec:method:benchmark}

Validation and evaluation require ground truth and instruments, released as
a two-sided benchmark \TODO{E1: extend to 1,300-book corpus}:

\textbf{Discovery questions.} Topical, known-item, and bibliographic-fact
questions, each with gold answers linked to catalog record identifiers so
correctness is verifiable by construction. Wording is LLM-polished into
natural patron language while gold answers remain frozen. Retrieval quality
is reported as nDCG@10~\cite{ndcg} with bootstrap confidence intervals and
paired permutation significance tests.

\textbf{Cataloging gold.} Records with subject-heading and classification
ground truth drawn from LC-derived and MARC-derived professional sources,
cross-checked against a second professional source (HathiTrust) to quantify
how much inter-source disagreement is inherent to the task
\TODO{E2: upgrade to panel-validated gold + agreement ceiling}. Because the
gold has multiple legitimate forms, acquisition quality is scored with a
multi-level rubric: \emph{exact}, \emph{semantic equivalence} (a different
valid heading for the same concept), and \emph{acceptable} (a correct
broader heading), plus classification top-$k$. Predicted headings that match
no gold are classified into an error taxonomy---over-specific,
under-specific, authority-violating (checked live against the LC authority
file), and valid-but-unmatched---so that ``wrong'' is not conflated with
``different but valid'' \TODO{E2: decompose with panel}.

\textbf{Relevance judgments.} Single-gold topical scoring is known to
understate systems that retrieve a relevant non-seed record, so we apply a
pooled graded relevance check: for each topical question the top records
from all configurations are pooled and a judge grades each candidate
(0/1/2) without knowing the seed; configurations are rescored with graded
nDCG@10. The LLM judge is calibrated against human raters
\TODO{E2: extend calibration to the full journal pool}.

\textbf{Generative-AI disclosure.} In line with the target venue's policy
\TODO{cite venue AI policy}, we disclose that \texttt{deepseek-chat} was
used to polish question wording, to generate experimental answers, and to
serve as the faithfulness judge; all generated content was programmatically
validated, and GenAI was not used to fabricate data or as an author.

\subsection{From measures to knowledge-engineering claims}
\label{sec:method:from-measures}

Three measured quantities support the paper's claims. \emph{Acquisition
accuracy} (Section~\ref{sec:valid}) is the multi-level agreement of acquired
knowledge with validated gold, read against an inter-cataloger agreement
ceiling rather than against exact match. \emph{Downstream utility}
(Section~\ref{sec:util}) is the retrieval-gain contribution of a knowledge
field, decomposed by representation (field ablation) and by knowledge-base
completeness (coverage sweep). \emph{Acquisition cost}
(Section~\ref{sec:cost}) is expert time to produce the same knowledge from
scratch versus verifying LLM output \TODO{E4}. Combining utility and cost
yields the deployment analysis (Section~\ref{sec:synth}): the value of
machine-acquired knowledge per expert-hour as a function of how complete the
knowledge base already is.
